\documentclass[11pt,a4paper]{article}
\usepackage[margin=25mm,headheight=15pt]{geometry}
\usepackage[T1]{fontenc}
\usepackage[utf8]{inputenc}
\usepackage{lmodern,amsmath,amssymb,amsthm,mathtools,bm}
\usepackage{xcolor,graphicx,booktabs,tabularx,array,longtable}
\usepackage{microtype,enumitem,fancyhdr}
\usepackage{xurl}
\usepackage{hyperref}
\definecolor{accent}{HTML}{163B55}
\hypersetup{colorlinks=true,linkcolor=accent,citecolor=accent,urlcolor=accent,
 pdftitle={Sharp fractional-power intertwining inequalities},pdfauthor={Matthew J. Colbrook},
 pdfsubject={AIM 897: mathematical result and proof}}
\urlstyle{same}
\setlength{\parindent}{0pt}
\setlength{\parskip}{0.35em}
\setlength{\emergencystretch}{3em}
\setlist{itemsep=0.25em,topsep=0.3em}
\newtheorem{theorem}{Theorem}[section]
\newtheorem{lemma}[theorem]{Lemma}
\newtheorem{proposition}[theorem]{Proposition}
\newtheorem{corollary}[theorem]{Corollary}
\theoremstyle{remark}
\newtheorem{remark}[theorem]{Remark}
\newtheorem{definition}[theorem]{Definition}
\DeclareMathOperator{\Tr}{Tr}
\DeclareMathOperator{\supp}{supp}
\DeclareMathOperator{\Lip}{Lip}
\DeclareMathOperator{\Rea}{Re}
\newcommand{\R}{\mathbb R}
\newcommand{\C}{\mathbb C}
\newcommand{\E}{\mathbb E}
\newcommand{\N}{\mathbb N}
\newcommand{\ip}[2]{\langle #1,#2\rangle}
\pagestyle{fancy}
\fancyhf{}
\fancyhead[L]{\small\sffamily AIM 897}
\fancyhead[R]{\small\sffamily Matthew J. Colbrook}
\fancyfoot[C]{\thepage}
\renewcommand{\headrulewidth}{0.3pt}
\allowdisplaybreaks[2]
\title{\textbf{Sharp fractional-power intertwining inequalities}}
\author{Matthew J. Colbrook\\[3pt]
{\small Department of Applied Mathematics and Theoretical Physics}\\
{\small University of Cambridge, United Kingdom}\\
{\small\href{mailto:mjc249@maths.cam.ac.uk}{\texttt{mjc249@maths.cam.ac.uk}}}}
\date{3 November 2026}
\begin{document}
\maketitle
\begin{abstract}
We prove the fractional-power intertwining inequality with sharp constant one for positive matrices and all exponents between zero and one. The argument yields spectral-floor refinements and a logarithmic commutator bound, and extends in operator norm to bounded operators and continuous nonnegative operator-monotone functions.
\end{abstract}
\noindent\textbf{Repository target.}
This manuscript concerns AIM problem 897 at revision
\texttt{5d7e26e8444e}. The precise scope is stated below.
\medskip

\section{Main result and proof}
\label{sec:897}

The original fractional-power question is recorded in \cite[Problem X.5.3]{Bhatia1997}.


\paragraph{Answer.} The constant-one inequality holds for every exponent
$0<r<1$. The main step is a sharp estimate for shifted square roots. A
positive integral representation proves the result above exponent one
half, and repeated square roots cover the remaining exponents.

\begin{theorem}\label{thm:897}
For every integer $n\ge1$, positive semidefinite matrices
$A,B\in M_n(\mathbb C)$, matrix $X\in M_n(\mathbb C)$, and $0<r<1$,
\[
 \boxed{\quad
 \|A^rX-XB^r\|
 \le\|X\|^{1-r}\|AX-XB\|^r.
 \quad}
\]
The norms are operator norms, and the constant one is optimal.
\end{theorem}

\subsection{An estimate at an extremal commutator vector}

\begin{lemma}\label{lem:897-orbit-energy}
Let $M\ge m_0I$ with $m_0\ge0$, and let $K^*=-K$ be finite complex
matrices. Put $k=\|K\|$, $H=[M,K]=MK-KM$, and $h=\|H\|>0$.
After replacing $K$ by $-K$ if necessary, choose a unit vector $w$ with
$Hw=hw$. Then
\[
 w^*Mw\ge m_0+\frac{h}{2k}.
 \tag{897.1}\label{eq:897-mean}
\]
\end{lemma}
\begin{proof}
The matrix $H$ is self-adjoint and its spectrum lies in $[-h,h]$.
A choice of sign ensures that $h$ is an eigenvalue. With $P=ww^*$,
\[
 H\ge h(2P-I).
 \tag{897.2}\label{eq:897-spectral-order}
\]
Indeed, both sides agree on the span of $w$, while on its orthogonal
complement the assertion is $H\ge-hI$. Also $k>0$, since $K=0$ would
imply $H=0$.

Set $U_t=e^{tK}$ for real $t$. These matrices are unitary. The matrix
$L=-iK$ is self-adjoint with spectrum in $[-k,k]$, and $U_t=e^{itL}$.
For $|t|\le\pi/(4k)$, spectral calculus gives
\[
 \operatorname{Re}(w^*U_tw)
 =w^*\cos(tL)w\ge\cos(k|t|)>0.
\]
It follows from \eqref{eq:897-spectral-order} that
\[
 \begin{split}
 g(t):=w^*U_t^*HU_tw
 &\ge h\bigl(2|w^*U_tw|^2-1\bigr)\\
 &\ge h\cos(2kt)
 \qquad\left(|t|\le\frac{\pi}{4k}\right).
 \end{split}
 \tag{897.3}\label{eq:897-orbit-bound}
\]
The function $f(t)=w^*U_t^*MU_tw$ satisfies $f(t)\ge m_0$ and
\[
 f'(t)=w^*U_t^*(MK-KM)U_tw=g(t).
\]
Integrating over the preceding quarter-period yields
\[
 \begin{split}
 w^*Mw=f(0)
 &=f\left(-\frac{\pi}{4k}\right)
   +\int_{-\pi/(4k)}^0g(t)\,dt\\
 &\ge m_0+h\int_{-\pi/(4k)}^0\cos(2kt)\,dt
 =m_0+\frac{h}{2k}.
 \end{split}
\]
\end{proof}

\begin{proposition}\label{prop:897-shifted-commutator}
For $M\ge m_0I\ge0$ and skew-adjoint $K$ with $k=\|K\|>0$, put
$d=\|[M^2,K]\|$. Then
\[
 \|[M,K]\|
 \le k\left(\sqrt{m_0^2+d/k}-m_0\right).
 \tag{897.4}\label{eq:897-shifted-commutator}
\]
\end{proposition}
\begin{proof}
If $h=\|[M,K]\|=0$, the assertion is immediate. Otherwise use the sign
choice and vector from Lemma~\ref{lem:897-orbit-energy}. With
$H=[M,K]$, the identity $[M^2,K]=MH+HM$ gives
\[
 d\ge\left|w^*[M^2,K]w\right|
   =2h\,w^*Mw
   \ge2m_0h+\frac{h^2}{k}.
\]
Solving this quadratic inequality for $h\ge0$ proves
\eqref{eq:897-shifted-commutator}. The sign change of $K$ leaves all
norms unchanged.
\end{proof}

\subsection{The shifted square-root inequality}

Let $A,B\ge0$, $X\ne0$, and write
\[
 k=\|X\|>0,\qquad d=\|AX-XB\|,\qquad\delta=d/k.
\]
For every $s\ge0$ we claim the sharper estimate
\[
 \|\sqrt{A+sI}\,X-X\sqrt{B+sI}\|
 \le k\bigl(\sqrt{s+\delta}-\sqrt s\bigr).
 \tag{897.5}\label{eq:897-shifted-root}
\]
To prove it, use the block matrices
\[
 M=\begin{pmatrix}\sqrt{A+sI}&0\\0&\sqrt{B+sI}\end{pmatrix},
 \qquad
 K=\begin{pmatrix}0&X\\-X^*&0\end{pmatrix}.
\]
Here $M\ge\sqrt s\,I$, $K^*=-K$, and $\|K\|=\|X\|=k$.
Writing $Y=\sqrt{A+sI}\,X-X\sqrt{B+sI}$ and $Z=AX-XB$ gives
\[
 [M,K]=\begin{pmatrix}0&Y\\Y^*&0\end{pmatrix},
 \qquad
 [M^2,K]=\begin{pmatrix}0&Z\\Z^*&0\end{pmatrix}.
\]
Their norms are $\|Y\|$ and $\|Z\|=d$, respectively; the terms $sX$
cancel in the second commutator. Proposition~\ref{prop:897-shifted-commutator}
with $m_0=\sqrt s$ proves \eqref{eq:897-shifted-root}.
At $s=0$ it gives the desired square-root inequality
\[
 \|A^{1/2}X-XB^{1/2}\|\le\|X\|^{1/2}\|AX-XB\|^{1/2}.
 \tag{897.6}\label{eq:897-root}
\]
In particular, if $d=0$, every shifted square-root commutator in
\eqref{eq:897-shifted-root} is zero.

\subsection{A positive integral for exponents above one half}

Fix $1/2<r<1$ and define
\[
 C_r^{-1}
 =\int_0^\infty\bigl(\sqrt{1+s}-\sqrt s\bigr)s^{r-3/2}\,ds.
 \tag{897.7}\label{eq:897-constant}
\]
This is a finite positive number. Near zero the integrand is bounded by
$s^{r-3/2}$, whose exponent is greater than $-1$. At infinity it is
bounded by $\tfrac12s^{r-2}$, whose exponent is less than $-1$.
Substituting $s=tu$ proves, for $t>0$,
\[
 t^r=C_r\int_0^\infty
       \bigl(\sqrt{t+s}-\sqrt s\bigr)s^{r-3/2}\,ds.
 \tag{897.8}\label{eq:897-integral}
\]
The same identity holds at $t=0$, where the integrand vanishes.
If an explicit special-function expression is desired, the identity
\[
 \sqrt{1+s}-\sqrt s=\frac12\int_0^1(s+u)^{-1/2}\,du
\]
and Tonelli's theorem give
\[
 C_r=\frac{2r}{\mathrm B(r-1/2,1-r)}
     =\frac{2r\sqrt\pi}{\Gamma(r-1/2)\Gamma(1-r)}.
\]
Only positivity and the normalization in \eqref{eq:897-constant} are
needed below.

Formula \eqref{eq:897-integral} applies to positive matrices by spectral
calculus, with convergence in operator norm. To verify the convergence
explicitly, for $A\ge0$ and $s>0$,
\[
 \|\sqrt{A+sI}-\sqrt s\,I\|
 =\sqrt{\|A\|+s}-\sqrt s
 \le\min\left\{\sqrt{\|A\|},\frac{\|A\|}{2\sqrt s}\right\}.
\]
The same integrable endpoint bounds used above therefore apply.
Consequently the scalar terms $\sqrt s\,X$ cancel to give
\[
 A^rX-XB^r
 =C_r\int_0^\infty
   \bigl(\sqrt{A+sI}\,X-X\sqrt{B+sI}\bigr)s^{r-3/2}\,ds.
\]
Using \eqref{eq:897-shifted-root}, positivity of the weight, and then
\eqref{eq:897-integral} with $t=\delta$ yields
\[
 \begin{split}
 \|A^rX-XB^r\|
 &\le kC_r\int_0^\infty
       \bigl(\sqrt{s+\delta}-\sqrt s\bigr)s^{r-3/2}\,ds\\
 &=k\delta^r
 =\|X\|^{1-r}\|AX-XB\|^r.
 \end{split}
 \tag{897.9}\label{eq:897-large-r}
\]
This also covers $d=0$, because both the integrand bound and $\delta^r$
are then zero.

\subsection{All remaining exponents and sharpness}

An inequality with constant one at exponents $\alpha$ and $\beta$ implies
it at their product. Indeed,
\[
 \begin{split}
 \|A^{\alpha\beta}X-XB^{\alpha\beta}\|
 &\le\|X\|^{1-\beta}\|A^\alpha X-XB^\alpha\|^\beta\\
 &\le\|X\|^{1-\alpha\beta}\|AX-XB\|^{\alpha\beta}.
 \end{split}
\]
For any $0<r\le1/2$, choose an integer $j\ge1$ such that
$\alpha=2^jr\in(1/2,1]$. The exponent $\alpha$ is covered by
\eqref{eq:897-large-r}, or is $1$, for which the inequality is an equality.
Compose it with $j$ applications of \eqref{eq:897-root} to obtain the
inequality at $r=\alpha2^{-j}$.

If $X=0$, both sides of Theorem~\ref{thm:897} are zero, so the case omitted
when dividing by $k$ is immediate. All $d=0$ cases were included above.
Finally, the one-dimensional choice $A=1$, $B=0$, $X=1$ gives equality for
every $r$, proving that the constant cannot be decreased.

\paragraph{Relation to the cited problems.}
Loring and Vides isolate the square-root constant-one target as
Conjecture~1.6 in their versioned preprint \cite{LoringVides2013}.
Herrera's later work \cite{Herrera2023} retains that question while proving
a factor below $1.00891$ for square roots and below $1.01975$ for the
broader family. The shifted estimate \eqref{eq:897-shifted-root} is the
step that makes the positive integral argument exact. The resulting
Theorem~\ref{thm:897} has all the quantifiers of the catalogue statement
\cite{AIM}.

\paragraph{Extension to bounded operators and unital $C^*$-algebras.}
The same inequalities hold in every unital $C^*$-algebra, with its
$C^*$-norm. To remove the finite-dimensional eigenvector assumption,
choose the sign of $H$ so that $h=\|H\|\in\sigma(H)$ and a state
$\varphi$ with $\varphi(H)=h$. In its GNS representation $(\pi,\xi)$,
$\pi(H)\xi=h\xi$: indeed, $0\le hI-H\le2hI$ implies
$\varphi((hI-H)^2)=0$. The orbit proof applies to $\pi(M),\pi(K),\xi$
using the bounds $\|\pi(H)\|\le h$ and $\|\pi(K)\|\le k$, and gives
$\varphi(M)\ge m_0+h/(2k)$. Thus
$\|[M^2,K]\|\ge|\varphi(MH+HM)|=2h\varphi(M)$, exactly as before.
The block matrices belong to $M_2(\mathcal A)$, and all functional-calculus
integrals converge in norm, so the rest of the proof is unchanged.

\subsection{A sharp bound with a positive spectral lower bound}

\begin{corollary}\label{cor:897-lower-bound}
Suppose $A,B\ge aI$ with $a\ge0$, and let $X\ne0$. Write
$k=\|X\|$ and $d=\|AX-XB\|$. Then, for every $0<r<1$,
\[
 \|A^rX-XB^r\|
 \le k\bigl((a+d/k)^r-a^r\bigr).
 \tag{897.10}\label{eq:897-positive-floor}
\]
If $a>0$, one also has
\[
 \|(\log A)X-X\log B\|
 \le k\log\left(1+\frac{d}{ak}\right).
 \tag{897.11}\label{eq:897-logarithm}
\]
Both bounds are sharp.
\end{corollary}
\begin{proof}
Put $\delta=d/k$. For the shifted square-root matrices used above, the
improved lower bound is $M\ge\sqrt{a+s}\,I$. Thus
Proposition~\ref{prop:897-shifted-commutator} gives
\[
 \|\sqrt{A+sI}\,X-X\sqrt{B+sI}\|
 \le k\bigl(\sqrt{a+s+\delta}-\sqrt{a+s}\bigr).
\]
For $1/2<r<1$, integrating this inequality with the same positive weight
as in \eqref{eq:897-integral} gives exactly
$k((a+\delta)^r-a^r)$. The case $r=1/2$ follows by setting $s=0$.

The refined estimate is also preserved under composition. If it holds at
$\alpha$, then $A^\alpha,B^\alpha\ge a^\alpha I$ and
\[
 \delta_\alpha:=k^{-1}\|A^\alpha X-XB^\alpha\|
 \le(a+\delta)^\alpha-a^\alpha.
\]
Applying the refined estimate at $\beta$ and using scalar monotonicity
therefore gives
\[
 k^{-1}\|A^{\alpha\beta}X-XB^{\alpha\beta}\|
 \le(a^\alpha+\delta_\alpha)^\beta-a^{\alpha\beta}
 \le(a+\delta)^{\alpha\beta}-a^{\alpha\beta}.
\]
The same square-root iteration used in the theorem proves
\eqref{eq:897-positive-floor} for all remaining $r$.

If $a>0$, divide \eqref{eq:897-positive-floor} by $r$ and let $r\downarrow0$.
Finite-dimensional spectral calculus gives convergence in operator norm
of the left-hand side to the logarithmic commutator, and the scalar
right-hand side tends to $k\log((a+\delta)/a)$. This proves
\eqref{eq:897-logarithm}. The one-dimensional choices
$A=a+\delta$, $B=a$, and $X=1$ give equality, with $a>0$ in the logarithmic
case.
\end{proof}

\subsection{A real-matrix reduction used during the investigation}

There is also a reduction of potential counterexamples to real matrices
of the same order, independent of the proof above. Diagonalize $A,B$ and set
$D_{ij}=a_i-b_j$ and $F_{ij}=a_i^r-b_j^r$. The two commutators become
$D\circ X$ and $F\circ X$, where $\circ$ denotes entrywise multiplication.
If $F\circ X=0$, there is no counterexample. Otherwise choose unit singular
vectors $u,v$ with $u^*(F\circ X)v=\|F\circ X\|$. Diagonal unitary
matrices $U,V$ can be chosen so that $u=U|u|$ and $v=V|v|$.
For $Z=U^*XV$, all three relevant norms are unchanged, and
\[
 |u|^{\mathsf T}(F\circ Z)|v|=\|F\circ X\|.
\]
Set $R=\operatorname{Re}Z=(Z+\overline Z)/2$. Conjugation preserves the
operator norm, and the multipliers $D,F$ are real. Hence
\[
 \|R\|\le\|X\|,\qquad
 \|D\circ R\|\le\|D\circ X\|,\qquad
 \|F\circ R\|\le\|F\circ X\|.
\]
The last inequality is an equality: taking real parts preserves the
displayed singular-vector witness. Thus $R$ preserves the numerator
and cannot increase either factor on the right. Any complex
counterexample would therefore yield a real one of the same order.

\subsection{Extension to nonnegative operator-monotone functions}
\label{sec:897-operator-monotone}

The shifted-square-root estimate also yields a broader consequence.
Here the norm remains the operator norm throughout.

\begin{theorem}\label{thm:897-operator-monotone}
Let $f:[0,\infty)\to[0,\infty)$ be continuous and operator monotone.
For positive elements $A,B$ of a unital $C^*$-algebra and any element $X$,
with $k=\|X\|>0$ and $\delta=\|AX-XB\|/k$, one has
\[
 \|f(A)X-Xf(B)\|\le k\bigl(f(\delta)-f(0)\bigr).
 \tag{897.12}\label{eq:897-monotone}
\]
The case $X=0$ has both sides zero. In particular, the weaker right-hand
side $k f(\delta)$ also holds. The scalar choices $A=\delta$, $B=0$,
and $X=1$ attain the displayed bound.
\end{theorem}
\begin{proof}
Call a continuous increasing function $g:[0,\infty)\to[0,\infty)$ with
$g(0)=0$ admissible if
$\|g(A)X-Xg(B)\|\le k g(\delta)$ for every such $A,B,X$.
Two admissible functions have admissible composition: applying the
outer estimate to the inner function and using scalar monotonicity proves
this directly. The class is also closed under nonnegative linear
combinations and limits that are uniform on compact intervals, by
continuous functional calculus and continuity of the norm.

For $h>0$, the function
\[
 g_h(t)=\frac{\sqrt{1+4ht}-1}{2h}
\]
is admissible by \eqref{eq:897-shifted-root}, applied with shift $1$ to
$4hA,4hB$. Fix $\tau>0$ and set $h=\tau/n$. If
$x_0=t>0$ and $x_j=g_h(x_{j-1})$, then
\[
 x_{j-1}=x_j+hx_j^2,\qquad 0<x_j\le t,
 \qquad
 \frac1{x_j}-\frac1{x_{j-1}}=\frac{h}{1+hx_j}.
\]
Summing the bounds $h/(1+ht)\le h/(1+hx_j)\le h$ gives
\[
 \frac{t}{1+\tau t}
 \le g_{\tau/n}^{\circ n}(t)
 \le\frac{t(1+\tau t/n)}{1+\tau t+\tau t/n}.
\]
The same bounds hold at $t=0$. The difference between the two endpoints is
\[
 \frac{\tau^2t^3}
 {n(1+\tau t)(1+\tau t+\tau t/n)}
 \le\frac{\tau^2R^3}{n}\qquad(0\le t\le R).
\]
Thus $t\mapsto t/(1+\tau t)$ is an admissible compact-uniform limit of
admissible iterates. Taking $\tau=1/s$ and multiplying by $1/s$ shows that
\[
 t\longmapsto\frac{t}{t+s}
 \quad\hbox{is admissible for every }s>0.
\]

The positive integral representation
\cite[Proposition~1.1]{Chansangiam2013}, absorbing its factor $1+s$ into
the measure on $(0,\infty)$, gives
\[
 f(t)-f(0)=bt+\int_{(0,\infty)}\frac{t}{t+s}\,d\mu(s),
 \qquad b\ge0,\quad
 \int_{(0,\infty)}\frac{d\mu(s)}{1+s}<\infty.
\]
For bounded positive operators the integral converges in norm: its integrand is bounded by
$\min\{1,\|A\|/s\}$, which is dominated by a constant times $(1+s)^{-1}$.
The identity function is admissible. Integrating the resolvent estimates
against the positive measure and using the triangle inequality therefore
proves \eqref{eq:897-monotone}; the scalar constant $f(0)X$ cancels.
\end{proof}

For a spectral floor $A,B\ge aI$, apply the theorem to
$t\mapsto f(a+t)-f(a)$ and to $A-aI,B-aI$. This gives the sharper
increment $k(f(a+\delta)-f(a))$ whenever $f$ is operator monotone on
$[a,\infty)$ and continuous at $a$. In particular it recovers the
fractional-power and logarithmic refinements above. No statement for
arbitrary unitarily invariant norms is being inferred from this argument.



\begin{thebibliography}{99}
\small
\setlength{\itemsep}{3pt}

\bibitem{LoringVides2013}
T.~A.~Loring and F.~Vides,
\emph{Estimating Norms of Commutators},
Experimental Math. \textbf{24} (2015), 106--122.
\href{https://arxiv.org/abs/1301.4252v3}{arXiv:1301.4252v3} (2013),
Conjecture 1.6 and Lemma 1.7.

\bibitem{Herrera2023}
D.~Herrera,
\emph{On the commutator modulus of continuity for operator monotone functions},
preprint, November 2023.
\href{https://arxiv.org/abs/2311.17448}{arXiv:2311.17448},
equation (7), Conjectures 1.8--1.9 and Theorems 1.11--1.12.

\bibitem{AIM}
MColbrook/AIM, \emph{AIM 897: Pedersen's sharp fractional-power commutator inequality},
repository statement at revision \texttt{5d7e26e8444e}.
\url{https://github.com/MColbrook/AIM/blob/5d7e26e8444eb98de91988454564e45e011107c4/problems/897-pedersen-s-sharp-fractional-power-commutator-inequality.md}.

\bibitem{Chansangiam2013}
P.~Chansangiam,
\emph{Integral Representations and Decompositions of Operator Monotone
Functions on the Nonnegative Reals}, preprint, 2013.
\href{https://arxiv.org/abs/1304.7936}{arXiv:1304.7936}, Proposition 1.1.

\bibitem{Bhatia1997}
R.~Bhatia, \emph{Matrix Analysis}, Graduate Texts in Mathematics
\textbf{169}, Springer, 1997, Chapter X, Problem X.5.3, p.~317,
and chapter notes, p.~321.
\href{https://doi.org/10.1007/978-1-4612-0653-8}{doi:10.1007/978-1-4612-0653-8}.

\end{thebibliography}
\end{document}
